\chapter{Financial Ratios, Benford's Law \& Audit Fraud Detection}

% ----------------------------------------------------------------------
% PART 1: EXECUTIVE ROADMAP & LEARNING OBJECTIVES
% ----------------------------------------------------------------------
\section{Executive Roadmap \& Learning Objectives}

Forensic ratio analysis and digital digit distribution analysis represent two foundational quantitative methodologies used to pierce corporate opacity. While financial ratios standardize operational results across scale and expose relational dislocations between operating efficiency, capital structure, and profitability, Benford's Law leverages statistical mechanics to identify anomalous human intervention in financial transaction journals.

\begin{tcolorbox}[colback=primarylight, colframe=primarydark, title={\textbf{Core Learning Competencies}}]
Upon mastering this chapter, students will be equipped to:
\begin{itemize}
    \item \textbf{Deconstruct Corporate Profitability via the DuPont Tree:} Dissect Return on Equity (ROE) into operating profit margin, total asset turnover, and the financial leverage multiplier.
    \item \textbf{Score Earnings Quality with the Beneish M-Score:} Formulate and evaluate the 8-variable probabilistic probit model to calculate financial manipulation probability.
    \item \textbf{Derive Benford's Law Invariants:} Prove the logarithmic probability distribution of leading digits from the principle of scale invariance.
    \item \textbf{Execute Statistical Audit Tests:} Apply Pearson's Chi-Square ($\chi^2$) goodness-of-fit and Mean Absolute Deviation (MAD) against empirical transaction ledgers.
    \item \textbf{Deploy Automated Digit Forensics:} Implement production-ready Python pipelines and Excel functions to flag non-conforming transaction clusters.
\end{itemize}
\end{tcolorbox}

\begin{table}[htbp]
\centering
\small
\begin{tabularx}{\textwidth}{llX}
\toprule
\textbf{Construct / Metric} & \textbf{Symbol} & \textbf{Mathematical Definition / Decision Rule} \\
\midrule
Return on Equity & $\text{ROE}$ & $\frac{\text{Net Income}}{\text{Shareholders' Equity}} = \left(\frac{NI}{\text{Sales}}\right) \times \left(\frac{\text{Sales}}{\text{Assets}}\right) \times \left(\frac{\text{Assets}}{\text{Equity}}\right)$ \\
Beneish M-Score & $M$ & $\boldsymbol{\beta}^T \mathbf{X} > -1.78 \implies \text{High Probability of Manipulation}$ \\
Days Sales Outstanding & $\text{DSO}$ & $\frac{\text{Average Accounts Receivable}}{\text{Net Credit Sales}} \times 365$ \\
Benford 1st-Digit Prob & $P(D_1 = d)$ & $\log_{10}\left(1 + \frac{1}{d}\right) \quad \text{for } d \in \{1, 2, \dots, 9\}$ \\
Benford 2nd-Digit Prob & $P(D_2 = d_2)$ & $\sum_{k=1}^9 \log_{10}\left(1 + \frac{1}{10k + d_2}\right) \quad \text{for } d_2 \in \{0, 1, \dots, 9\}$ \\
Goodness-of-Fit Statistic & $\chi^2$ & $\sum_{d=1}^9 \frac{(O_d - E_d)^2}{E_d} \quad (df = 8)$ \\
Mean Absolute Deviation & $\text{MAD}$ & $\frac{1}{9} \sum_{d=1}^9 |p_d^{\text{obs}} - p_d^{\text{exp}}|$ \\
\bottomrule
\end{tabularx}
\caption{Mathematical Notation Guide for Ratio Decomposition and Benford Forensics.}
\label{tab:w5_notation}
\end{table}

% ----------------------------------------------------------------------
% PART 2: MATHEMATICAL & COMPUTATIONAL FOUNDATIONS
% ----------------------------------------------------------------------
\section{Mathematical \& Computational Foundations}

\subsection{The DuPont 3-Step and 5-Step Decompositions}
Developed by Donaldson Brown at the DuPont Corporation in 1914, the DuPont model separates an enterprise's return on equity into operational, capital efficiency, and leverage dimensions:
\begin{equation}
\text{ROE} = \underbrace{\left(\frac{\text{Net Income}}{\text{Revenue}}\right)}_{\text{Net Profit Margin}} \times \underbrace{\left(\frac{\text{Revenue}}{\text{Total Assets}}\right)}_{\text{Asset Turnover}} \times \underbrace{\left(\frac{\text{Total Assets}}{\text{Shareholders' Equity}}\right)}_{\text{Equity Multiplier}}
\label{eq:dupont_3step}
\end{equation}
The extended 5-step model further isolates operating profit from financing interest burden and corporate taxation:
\begin{equation}
\text{ROE} = \underbrace{\left(\frac{\text{Net Income}}{\text{EBT}}\right)}_{\text{Tax Burden}} \times \underbrace{\left(\frac{\text{EBT}}{\text{EBIT}}\right)}_{\text{Interest Burden}} \times \underbrace{\left(\frac{\text{EBIT}}{\text{Revenue}}\right)}_{\text{Operating Margin}} \times \left(\frac{\text{Revenue}}{\text{Assets}}\right) \times \left(\frac{\text{Assets}}{\text{Equity}}\right)
\label{eq:dupont_5step}
\end{equation}

\subsection{The Beneish M-Score Probabilistic Model}
Formulated by Professor Messod Beneish in 1999, the 8-variable M-Score is a probit model calibrated on known accounting restatement and financial manipulation firms:
\begin{align}
M &= -4.84 + 0.920 \cdot \text{DSRI} + 0.528 \cdot \text{GMI} + 0.404 \cdot \text{AQI} + 0.892 \cdot \text{SGI} \nonumber \\
&\quad + 0.115 \cdot \text{DEPI} - 0.172 \cdot \text{SGAI} + 4.037 \cdot \text{TATA} + 0.0327 \cdot \text{LVGI}
\label{eq:beneish_m}
\end{align}
where:
\begin{itemize}
    \item $\text{DSRI} = \frac{\text{Receivables}_t / \text{Sales}_t}{\text{Receivables}_{t-1} / \text{Sales}_{t-1}}$ (Days Sales in Receivables Index)
    \item $\text{GMI} = \frac{\text{Gross Margin}_{t-1}}{\text{Gross Margin}_t}$ (Gross Margin Deterioration Index)
    \item $\text{AQI} = \frac{1 - (\text{CA}_t + \text{PP\&E}_t + \text{Sec}_t)/\text{Assets}_t}{1 - (\text{CA}_{t-1} + \text{PP\&E}_{t-1} + \text{Sec}_{t-1})/\text{Assets}_{t-1}}$ (Asset Quality Index)
    \item $\text{SGI} = \frac{\text{Sales}_t}{\text{Sales}_{t-1}}$ (Sales Growth Index)
    \item $\text{DEPI} = \frac{\text{DeprRate}_{t-1}}{\text{DeprRate}_t}$ (Depreciation Slowdown Index)
    \item $\text{SGAI} = \frac{\text{SG\&A}_t / \text{Sales}_t}{\text{SG\&A}_{t-1} / \text{Sales}_{t-1}}$ (Sales Overhead Expense Index)
    \item $\text{LVGI} = \frac{\text{Leverage}_t}{\text{Leverage}_{t-1}}$ (Leverage Index)
    \item $\text{TATA} = \frac{\text{Operating Income}_t - \text{CFO}_t}{\text{Total Assets}_t}$ (Total Accruals to Total Assets)
\end{itemize}
\textbf{Decision Threshold:} If $M > -1.78$, the probability of manipulation is statistically elevated ($p > 3.8\%$). In Enron's 2000 financial filings, the Beneish M-Score reached $-1.49$, flagging extensive manipulation prior to bankruptcy.

\subsection{Benford's Law: First and Second-Digit Distributions}
Discovered by Simon Newcomb (1881) and popularized by Frank Benford (1938), Benford's Law describes the natural frequency distribution of leading digits in multi-scale numerical phenomena:
\begin{equation}
P(D_1 = d) = \log_{10}\left(1 + \frac{1}{d}\right) = \log_{10}(d + 1) - \log_{10}(d), \quad d \in \{1, 2, \dots, 9\}
\label{eq:benford_1st}
\end{equation}
The second-digit distribution expands to:
\begin{equation}
P(D_2 = d_2) = \sum_{k=1}^9 \log_{10}\left(1 + \frac{1}{10k + d_2}\right), \quad d_2 \in \{0, 1, \dots, 9\}
\label{eq:benford_2nd}
\end{equation}

% ----------------------------------------------------------------------
% PART 3: VISUAL & GEOMETRIC INTUITION
% ----------------------------------------------------------------------
\section{Visual \& Geometric Intuition: The DuPont Decomposition Tree}

Figure~\ref{fig:dupont_tree} displays the structural hierarchy of Return on Equity, demonstrating how operational profit margins, asset turnover, and balance sheet leverage multiply to determine total equity yield.

\begin{figure}[htbp]
\centering
\begin{tikzpicture}[
    font=\sffamily\footnotesize,
    node distance=1.2cm and 1.6cm,
    head/.style={rectangle, draw=primarydark, fill=primarylight, thick, rounded corners=4pt, minimum width=3.4cm, minimum height=1.1cm, align=center},
    dstep/.style={rectangle, draw=tealborder, fill=tealbg, thick, rounded corners=3pt, minimum width=3.2cm, minimum height=1.0cm, align=center},
    leaf/.style={rectangle, draw=coralborder, fill=coralbg, thick, rounded corners=3pt, minimum width=3.0cm, minimum height=0.9cm, align=center},
    arr/.style={->, >=stealth, thick, color=primarydark}
]

% Root
\node[head] (roe) {\textbf{Return on Equity (ROE)}\\$\text{ROE} = \frac{\text{Net Income}}{\text{Shareholders' Equity}}$};

% Level 1
\node[dstep, below left=1.2cm and 1.5cm of roe] (npm) {\textbf{Net Profit Margin}\\$\frac{\text{Net Income}}{\text{Sales}}$\\(Operational Margin)};
\node[dstep, below=1.2cm of roe] (ato) {\textbf{Total Asset Turnover}\\$\frac{\text{Sales}}{\text{Total Assets}}$\\(Asset Utilization)};
\node[dstep, below right=1.2cm and 1.5cm of roe] (em) {\textbf{Equity Multiplier}\\$\frac{\text{Total Assets}}{\text{Equity}}$\\(Financial Leverage)};

% Level 2 Drivers
\node[leaf, below=0.8cm of npm] (ni_driver) {Net Income Drivers\\$GP - \text{SG\&A} - \text{Taxes}$};
\node[leaf, below=0.8cm of ato] (inv_driver) {Asset Drivers\\$CA + NCA + \text{Inv}$};
\node[leaf, below=0.8cm of em] (debt_driver) {Capital Structure\\$\text{Debt} / \text{Equity}$};

% Connectors
\draw[arr] (roe) -- (npm);
\draw[arr] (roe) -- (ato);
\draw[arr] (roe) -- (em);
\draw[arr] (npm) -- (ni_driver);
\draw[arr] (ato) -- (inv_driver);
\draw[arr] (em) -- (debt_driver);

\end{tikzpicture}
\caption{The DuPont Three-Step Profitability Decomposition Tree.}
\label{fig:dupont_tree}
\end{figure}

% ----------------------------------------------------------------------
% PART 4: DEEP INTUITION & CONCEPTUAL DERIVATIONS
% ----------------------------------------------------------------------
\section{Deep Intuition \& Conceptual Derivations}

\subsection{Why Benford's Law Emerges from Scale Invariance}
Why should the digit $1$ appear as the leading digit $30.1\%$ of the time, while digit $9$ appears only $4.6\%$ of the time?

Consider a global currency exchange. If an empirical financial transaction dataset conforms to a universal statistical distribution, the relative probability of leading digits cannot depend on whether transactions are measured in Indian Rupees (INR), US Dollars (USD), or Euros (EUR). That is, the distribution must be \index{scale invariance}\textbf{scale-invariant}:
\begin{equation}
f(\lambda x) \, d(\lambda x) = f(x) \, dx \implies f(\lambda x) \lambda = f(x)
\end{equation}
Setting $x = 1$ and letting $\lambda$ vary yields the functional equation:
\begin{equation}
f(\lambda) = \frac{f(1)}{\lambda} \propto \frac{1}{\lambda}
\end{equation}
Thus, the underlying probability density function must be reciprocal: $f(x) \propto \frac{1}{x}$. Integrating this continuous density over the interval where the first digit equals $d$ within any decade $[10^k, 10^{k+1})$:
\begin{align}
P(D_1 = d) &= \int_{d \cdot 10^k}^{(d+1) \cdot 10^k} \frac{C}{x} \, dx = C \left[ \ln((d+1) \cdot 10^k) - \ln(d \cdot 10^k) \right] \nonumber \\
&= C \ln\left(\frac{d+1}{d}\right) = C \ln\left(1 + \frac{1}{d}\right)
\end{align}
Since $\sum_{d=1}^9 \log_{10}\left(1 + \frac{1}{d}\right) = \log_{10}\left(\prod_{d=1}^9 \frac{d+1}{d}\right) = \log_{10}(10) = 1$, setting base-$10$ logarithms naturally normalizes the distribution!

% ----------------------------------------------------------------------
% PART 5: EXPLICIT GOTCHAS & DIAGNOSTIC MISCONCEPTIONS
% ----------------------------------------------------------------------
\section{Explicit Gotchas \& Diagnostic Misconceptions}

\begin{trapbox}{Gotcha 1: Applying Benford's Law to Constrained Datasets}
\textbf{The Trap:} Running Benford's Law on employee salary bands, flight numbers, pin codes, or retail prices fixed at Rs.~$99$ or Rs.~$499$.\\
\textbf{Forensic Reality:} Benford's Law applies only to data that spans multiple orders of magnitude, is unconstrained by regulatory caps or artificial rounding conventions, and results from multiplicative processes. Applying it to human heights, ZIP codes, or pre-set pricing bands produces severe false-positive fraud flags.
\end{trapbox}

\begin{trapbox}{Gotcha 2: The Large Sample $\chi^2$ P-Value Trap}
\textbf{The Trap:} Concluding that an ERP database of $N = 2{,}000{,}000$ transactions is fraudulent simply because the Chi-Square test yields $p < 0.0001$.\\
\textbf{Forensic Reality:} The Pearson $\chi^2$ statistic is linearly proportional to sample size: $\chi^2 = N \sum \frac{(p_{\text{obs}} - p_{\text{exp}})^2}{p_{\text{exp}}}$. For very large $N$, infinitesimal, commercially benign deviations reject the null hypothesis. Auditors must rely on \textbf{Mean Absolute Deviation (MAD)} or \textbf{Z-score tests per digit} rather than aggregate $\chi^2$ $p$-values.
\end{trapbox}

\begin{trapbox}{Gotcha 3: The Unadjusted DuPont Leverage Distortion}
\textbf{The Trap:} Assuming that a company whose Return on Equity expanded from $14\%$ to $26\%$ has achieved brilliant operational improvement.\\
\textbf{Forensic Reality:} Under the DuPont equation, $\text{ROE} = \text{NPM} \times \text{ATO} \times \text{EM}$. If operating profit margin and asset turnover both halved while management quadrupled debt to repurchase equity, the Equity Multiplier spikes, creating an artificial surge in ROE while pushing the firm to the brink of bankruptcy.
\end{trapbox}

\begin{trapbox}{Gotcha 4: Misinterpreting Beneish M-Score Thresholds}
\textbf{The Trap:} Treating an M-Score of $-1.50$ ($> -1.78$) as definitive proof that the CFO has committed criminal fraud.\\
\textbf{Forensic Reality:} The Beneish M-Score is a statistical screening instrument indicating elevated risk, not legal proof of fraud. Hyper-growth tech companies with rapidly expanding receivables and heavy capital investments often generate high M-Scores without fraudulent accounting.
\end{trapbox}

% ----------------------------------------------------------------------
% PART 6: COMPARATIVE SYNTHESIS & SUMMARY TABLE
% ----------------------------------------------------------------------
\section{Comparative Synthesis \& Summary Table}

Table~\ref{tab:comparison_benford_mad} outlines Nigrini's authoritative Mean Absolute Deviation (MAD) conformity benchmarks for first-digit Benford audit tests.

\begin{table}[htbp]
\centering
\small
\begin{tabularx}{\textwidth}{llXX}
\toprule
\textbf{Conformity Level} & \textbf{First-Digit MAD Range} & \textbf{Audit Interpretation} & \textbf{Recommended Forensic Action} \\
\midrule
\textbf{Close Conformity} & $\text{MAD} \le 0.006$ & Data closely conforms to natural multiplicative distribution. & Standard procedural testing; no anomalous digit clustering. \\
\addlinespace
\textbf{Acceptable Conformity} & $0.006 < \text{MAD} \le 0.012$ & Minor statistical variations within normal enterprise operating bounds. & No immediate red flags; routine audit procedures. \\
\addlinespace
\textbf{Marginally Acceptable} & $0.012 < \text{MAD} \le 0.015$ & Modest digit deviations; potential threshold manipulation. & Expand sample size; examine specific spikes in digits $4, 5, 9$. \\
\addlinespace
\textbf{Non-Conforming} & $\text{MAD} > 0.015$ & Severe distortion of natural leading digit frequencies. & High audit risk; initiate targeted forensic journal vouchers audit. \\
\bottomrule
\end{tabularx}
\caption{Nigrini's Mean Absolute Deviation (MAD) Compliance Benchmarks for First-Digit Audit Testing.}
\label{tab:comparison_benford_mad}
\end{table}

% ----------------------------------------------------------------------
% PART 7: PYTHON / SCIENTIFIC COMPUTING LAB
% ----------------------------------------------------------------------
\section{Python / Scientific Computing Lab: Automated Benford Testing Engine}

Listing~\ref{lst:benford_engine} implements an automated audit engine that computes empirical first-digit distributions, compares them against theoretical Benford probabilities, calculates $\chi^2$ and MAD metrics, and flags high-risk transaction clusters.

\begin{lstlisting}[language=Python, caption={Benford First-Digit Audit Engine with $\chi^2$ and MAD in Python}, label={lst:benford_engine}]
import numpy as np
import pandas as pd
from scipy.stats import chisquare

def audit_benford_first_digit(amounts: np.ndarray, alpha: float = 0.01) -> dict:
    """
    Performs Benford's Law First-Digit Analysis on financial transaction amounts.
    
    Parameters:
        amounts (np.ndarray): Array of financial transaction magnitudes (positive floats).
        alpha (float): Significance level for Chi-Square goodness-of-fit test.
        
    Returns:
        dict: Complete audit report containing observed/expected tables, Chi2, and MAD.
    """
    # 1. Clean data: Filter non-positive values
    valid_amounts = amounts[amounts > 0]
    n = len(valid_amounts)
    if n < 100:
        raise ValueError("Benford audit requires at least 100 observations for statistical validity.")
        
    # 2. Extract first digit using string conversion
    first_digits = np.array([int(str(f"{val:.8e}")[0]) for val in valid_amounts])
    
    # 3. Calculate observed counts and theoretical Benford probabilities
    digits = np.arange(1, 10)
    observed_counts = np.array([np.sum(first_digits == d) for d in digits])
    observed_props = observed_counts / n
    expected_props = np.log10(1.0 + 1.0 / digits)
    expected_counts = expected_props * n
    
    # 4. Compute Test Statistics
    chi2_stat, p_val = chisquare(observed_counts, f_exp=expected_counts)
    mad = np.mean(np.abs(observed_props - expected_props))
    
    # 5. Evaluate Nigrini MAD Compliance
    if mad <= 0.006:
        conformity = "Close Conformity"
    elif mad <= 0.012:
        conformity = "Acceptable Conformity"
    elif mad <= 0.015:
        conformity = "Marginally Acceptable"
    else:
        conformity = "Non-Conforming (High Audit Risk)"
        
    results_df = pd.DataFrame({
        'Digit': digits,
        'Observed_Count': observed_counts,
        'Expected_Count': expected_counts.round(1),
        'Observed_Pct': (observed_props * 100).round(2),
        'Expected_Pct': (expected_props * 100).round(2),
        'Difference_Pct': ((observed_props - expected_props) * 100).round(2)
    })
    
    return {
        'sample_size': n,
        'chi2_statistic': round(chi2_stat, 3),
        'p_value': p_val,
        'mad': round(mad, 5),
        'conformity_status': conformity,
        'flagged': (p_val < alpha) or (mad > 0.015),
        'digit_distribution': results_df
    }

# -------------------------------------------------------------
# Test Dataset with Controlled Approval-Bypassing Cluster
# -------------------------------------------------------------
if __name__ == '__main__':
    np.random.seed(42)
    # Natural log-normal commercial invoice volume
    natural_txns = np.random.lognormal(mean=8.5, sigma=1.8, size=8000)
    # Fraudulent invoice cluster just below the Rs. 50,000 dual-signoff limit
    fraud_cluster = np.random.uniform(48000, 49999, size=1200)
    
    ledger_data = np.concatenate([natural_txns, fraud_cluster])
    audit = audit_benford_first_digit(ledger_data)
    
    print(f"Total Transactions: {audit['sample_size']}")
    print(f"Chi-Square: {audit['chi2_statistic']}, p-val: {audit['p_value']:.4e}")
    print(f"MAD: {audit['mad']} -> {audit['conformity_status']}")
    print(f"Audit Status: {'FLAGGED FOR INVESTIGATION' if audit['flagged'] else 'CLEAN'}\n")
    print(audit['digit_distribution'].to_string(index=False))
\end{lstlisting}

% ----------------------------------------------------------------------
% PART 8: COMPREHENSIVE PRACTICE PROBLEMS
% ----------------------------------------------------------------------
\section{Comprehensive Practice Problems}
\textit{Complete mathematical derivations and diagnostic MCQ distractor rationales are located in the Appendix.}

\subsection*{Part A: Analytical \& Computational Problems}

\begin{enumerate}[label=\textbf{P5.\arabic*}]
    \item \textbf{Benford Distribution Probability Proof \& Digit Calculations:}\\
    Using the foundational Benford formulation $P(D_1 = d) = \log_{10}(1 + 1/d)$:
    \begin{enumerate}[label=(\alph*)]
        \item Compute exact theoretical probabilities for $d = 1$, $d = 3$, and $d = 8$ to four decimal places.
        \item Prove analytically that the sum of probabilities across all nine digits $\sum_{d=1}^9 P(D_1 = d) \equiv 1.0$.
        \item Formulate the second-digit formula $P(D_2 = d_2)$ and compute the probability that $D_2 = 0$.
    \end{enumerate}

    \item \textbf{Substantive Chi-Square \& MAD Goodness-of-Fit Test:}\\
    An internal audit department extracts $N = 10{,}000$ expense vouchers submitted by an executive sales division. The observed leading digit counts for digits $1$ through $9$ are:
    \begin{equation*}
    [2100, 1400, 1100, 1900, 1300, 800, 600, 450, 350]
    \end{equation*}
    \begin{enumerate}[label=(\alph*)]
        \item Calculate the expected theoretical Benford counts for each digit $d \in \{1, \dots, 9\}$.
        \item Compute the Pearson Chi-Square test statistic $\chi^2$ and test the null hypothesis of conformity at $\alpha = 0.01$ ($df = 8$, critical value $\chi^2_{0.01} = 20.09$).
        \item Calculate the Mean Absolute Deviation (MAD) and evaluate Nigrini's compliance rating.
        \item Which digit accounts for the greatest anomaly, and what forensic hypothesis explains this spike?
    \end{enumerate}

    \item \textbf{DuPont Decomposition \& Capital Structure Arbitrage:}\\
    Company $A$ and Company $B$ both report a Return on Equity (ROE) of $24\%$ in FY2025. Financial records reveal:
    \begin{center}
    \begin{tabular}{lrr}
    \toprule
    \textbf{Metric} & \textbf{Company A} & \textbf{Company B} \\
    \midrule
    Net Sales Revenue & Rs.~$1{,}000$ Crores & Rs.~$400$ Crores \\
    Net Income & Rs.~$120$ Crores & Rs.~$24$ Crores \\
    Total Assets & Rs.~$600$ Crores & Rs.~$800$ Crores \\
    Shareholders' Equity & Rs.~$500$ Crores & Rs.~$100$ Crores \\
    \bottomrule
    \end{tabular}
    \end{center}
    \begin{enumerate}[label=(\alph*)]
        \item Deconstruct the ROE of both firms using the 3-Step DuPont model.
        \item Contrast the financial risk, operational margin, and earnings sustainability of the two firms.
    \end{enumerate}

    \item \textbf{Beneish M-Score Earnings Quality Assessment:}\\
    A healthcare provider reports Days Sales in Receivables Index ($\text{DSRI} = 1.62$), Gross Margin Index ($\text{GMI} = 1.35$), Asset Quality Index ($\text{AQI} = 1.28$), Sales Growth Index ($\text{SGI} = 1.45$), Depreciation Index ($\text{DEPI} = 1.12$), SG\&A Expense Index ($\text{SGAI} = 0.95$), Total Accruals to Total Assets ($\text{TATA} = 0.09$), and Leverage Index ($\text{LVGI} = 1.25$).
    \begin{enumerate}[label=(\alph*)]
        \item Calculate the Beneish 8-variable M-Score.
        \item Interpret the result against the $-1.78$ threshold and identify the three most aggressive variables.
    \end{enumerate}

    \item \textbf{Receivables Aging Latency \& Days Sales Outstanding (DSO):}\\
    A wholesale distributor reports credit sales of Rs.~$73$ Crores and ending accounts receivable of Rs.~$24$ Crores in Year 3. In Year 1, credit sales were Rs.~$60$ Crores with ending receivables of Rs.~$8$ Crores.
    \begin{enumerate}[label=(\alph*)]
        \item Compute DSO for Year 1 and Year 3 assuming 365 days.
        \item Formulate three forensic hypotheses explaining why DSO expanded threefold.
    \end{enumerate}
\end{enumerate}

\subsection*{Part B: Diagnostic Multiple-Choice Questions}

\begin{enumerate}[label=\textbf{MCQ 5.\arabic*}]
    \item When auditing an enterprise general ledger containing $1{,}500{,}000$ procurement transactions, an auditor computes the Pearson Chi-Square statistic against Benford's Law, obtaining $\chi^2 = 34.2$ with $p = 0.00002$. However, the calculated Mean Absolute Deviation (MAD) is $0.0038$. What is the correct forensic interpretation?
    \begin{enumerate}[label=(\alph*)]
        \item The procurement ledger is severely fraudulent because the $p$-value is far below $0.01$.
        \item The rejection of the null hypothesis is an artifact of the very large sample size; the low MAD confirms close conformity with natural Benford distribution.
        \item The dataset must be discarded because Benford's Law is statistically invalid for datasets over $50{,}000$ transactions.
        \item The auditor should immediately refer the procurement team for criminal indictment based on the low $p$-value.
    \end{enumerate}

    \item In the 8-variable Beneish M-Score model, which of the following metrics carries the single highest positive regression weight, meaning an increase in this variable most aggressively raises the probability of financial manipulation?
    \begin{enumerate}[label=(\alph*)]
        \item Total Accruals to Total Assets (TATA), with a weight of $+4.037$.
        \item Sales Growth Index (SGI), with a weight of $+0.892$.
        \item Asset Quality Index (AQI), with a weight of $+0.404$.
        \item Gross Margin Index (GMI), with a weight of $+0.528$.
    \end{enumerate}
\end{enumerate}
